\documentclass{article}
\usepackage[utf8]{inputenc}
\usepackage{amsmath}
\usepackage{amssymb}
\usepackage{geometry}
\geometry{a4paper, margin=1in}
\title{Reconstruction of Neural Tangent Kernel (NTK) Calculations for an MMNN Network}
\author{Analysis of provided notes}
\date{\today}

\begin{document}

\maketitle

\section{Introduction and Definitions}
This document first introduces the Multi-component and Multi-layer Neural Network (MMNN) architecture in its general n-dimensional form. It then reconstructs a series of calculations aimed at deriving the neural tangent kernel (NTK) for this network. The central hypothesis for the NTK calculation, in accordance with learning strategy S1 from the source paper, is that only the parameters of the external linear combination layer are trained, while the weights and biases internal to the activation function are fixed and randomly initialized.

\subsection{MMNN Architecture}
A single layer of an MMNN is a shallow neural network approximating a vector-valued function $\bmf:\mathbb{R}^{d_{in}} \to \mathbb{R}^{d_{out}}$. It is defined as:
\begin{equation}
    \bmh(\bmx) = \bmA \sigma(\bmW\bmx + \bmb) + \bmc
\end{equation}
This can be written in a more compact form to better illustrate its structure:
\begin{equation}
    \bmh(\bmx) = \tilde{\bm{A}}
    \begin{pmatrix}
        \sigma(\tilde{\bm{W}}\tilde{\bm{x}}) \\
        1
    \end{pmatrix}
\end{equation}
where $\tilde{\bm{W}} = [\bm{W}, \bm{b}]$, $\tilde{\bm{A}} = [\bm{A}, \bm{c}]$, and $\tilde{\bm{x}} = [\bm{x}^{\top}, 1]^{\top}$.

We call each element of $\bmh$, i.e., $\bmh[i] = \tilde{\bm{A}}[i, :] \cdot ([\sigma(\tilde{\bm{W}}\tilde{\bm{x}})]^{\top}, 1)^{\top}$, a \textit{component}. This representation highlights several key features:
\begin{enumerate}
    \item Each component is a linear combination of a shared set of basis functions, $\{\sigma((\bm{W}\bm{x}+\bm{b})_j)\}_{j=1}^n$.
    \item Different components of $\bmh$ use the same basis functions but with different linear coefficients, given by the rows of $\tilde{\bm{A}}$.
    \item According to learning strategy S1, only the external coefficients $\tilde{\bm{A}}$ are trained, while the internal parameters $\tilde{\bm{W}}$ that define the basis are randomly initialized and fixed.
    \item The output dimension $d_{out}$ (the number of components, or rank) and the network width $n$ (the number of basis functions) are tunable hyperparameters. Typically, $d_{out} \ll n$.
\end{enumerate}

In contrast to a typical FCNN layer, which takes the form $\sigma(\tilde{\bm{W}}\tilde{\bm{x}})$ and in which all weights are trained, an MMNN layer is composed of multiple components where each is a linear combination of randomly parameterized hidden neurons. This structure acts as a smooth decomposition that can be trained more effectively.

The decomposition theorems presented earlier provide a mathematical intuition for this architecture. The fixed, random basis functions generated by $\sigma(\tilde{\bm{W}}\tilde{\bm{x}})$ play a role analogous to the piecewise linear transformations $\psi_i$. The trainable parameters $\tilde{\bm{A}}$ then learn the appropriate linear combination of these basis functions, conceptually similar to learning the smoother functions $f_i$ in the decomposition. This parallel between the parameters $(\tilde{\bm{A}}, \tilde{\bm{W}})$ and the decomposition functions $(f_i, \psi_i)$ will be explored further in the report.

An MMNN is a full multi-layer composition of such blocks:
\begin{equation}
    \bmh_{total} = \bmh_m \circ \cdots \circ \bmh_2 \circ \bmh_1
\end{equation}
where each block $\bmh_i: \mathbb{R}^{d_{i-1}} \to \mathbb{R}^{d_i}$ is a multi-component shallow network. The overall network is characterized by its width (largest number of basis functions), rank (largest intermediate dimension), and depth (number of layers).

\subsection{NTK Preliminaries}
For the NTK calculation, we consider the trainable parameters $\theta = \{A, c\}$. The kernel is defined by the dot product of the network output gradients with respect to these parameters:
\begin{equation}
    K_{NTK}(x, x') = \mathbb{E}_{\theta_{init}} \left[ \left\langle \nabla_{\theta} h(x), \nabla_{\theta} h(x') \right\rangle \right]
\end{equation}
For a single layer, the gradient with respect to the coefficients of $A$ is the vector of activations $\sigma(Wx+b)$, and the gradient with respect to $c$ is a vector of 1s.

\section{Kernel Calculation for One Layer}
The kernel calculation reduces to the expectation of the product of activation functions for two different inputs, $x$ and $x'$.

\subsection{Problem Reduction}
Let two Gaussian variables $g_1$ and $g_2$ be defined by:
\begin{align*}
    g_1 &= w^T x + b \\
    g_2 &= w^T x' + b
\end{align*}
where the weights $w$ and bias $b$ are assumed to follow a standard normal distribution, $w \sim \mathcal{N}(0, I)$ and $b \sim \mathcal{N}(0, \sigma_b^2)$. The expectation to calculate is:
\begin{equation}
    \mathbb{E}_{w,b} [\sigma(g_1) \sigma(g_2)]
\end{equation}
The pair $(g_1, g_2)$ follows a bivariate normal distribution with mean $(\mu, \mu)$ and covariance matrix $\Sigma$. If we assume that $x$ and $x'$ are normalized ($\|x\|=\|x'\|=1$) and $b$ is centered ($\mathbb{E}[b]=0$), then the variables are centered and their covariance is:
\begin{equation}
    \Sigma = \begin{pmatrix}
        \|x\|^2 + \sigma_b^2 & \langle x, x' \rangle + \sigma_b^2 \\
        \langle x, x' \rangle + \sigma_b^2 & \|x'\|^2 + \sigma_b^2
    \end{pmatrix}
\end{equation}
The correlation coefficient is $\rho = \langle x, x' \rangle$.

In the simplified case where $\sigma_b=0$ and $\|x\|=\|x'\|=1$, the calculation reduces to:
\begin{equation}
    I = \mathbb{E}[\sigma(Z_1)\sigma(Z_2)]
\end{equation}
where $(Z_1, Z_2)$ is a standard bivariate normal distribution (mean 0, variance 1) with correlation $\rho$. For ReLU activation, $\sigma(z) = \max(0,z)$, the result is known:
\begin{equation}
    \mathbb{E}[\text{ReLU}(Z_1)\text{ReLU}(Z_2)] = \frac{1}{2\pi} (\sqrt{1-\rho^2} + (\pi - \arccos\rho)\rho)
\end{equation}

\subsection{General Non-Centered Case}
For the more general case where the Gaussian variables are not centered, $g \sim \mathcal{N}(\mu, \sigma^2)$ and $g' \sim \mathcal{N}(\mu, \sigma^2)$ with correlation $\rho$, the expectation $I = \mathbb{E}[\sigma(g)\sigma(g')]$ is calculated as follows.

Let $a = -\frac{\mu}{\sigma}$. The expectation can be decomposed into:
\begin{equation}
    I = \sigma^2 \mathbb{E}[Z Z' \cdot \mathbf{1}_{Z>a, Z'>a}] + 2\mu\sigma \mathbb{E}[Z \cdot \mathbf{1}_{Z>a, Z'>a}] + \mu^2 \mathbb{E}[\mathbf{1}_{Z>a, Z'>a}]
\end{equation}
where $Z, Z'$ are standard normal variables. This decomposes into three integrals:
\begin{equation}
    I = \sigma^2 I_2 + 2\mu\sigma I_1 + \mu^2 I_0
\end{equation}

The formulas for these terms (marked as coming from the bibliography in the notes) are:
\begin{itemize}
    \item \textbf{Term $I_0$ (Probability)}:
    \begin{equation}
        I_0 = P(Z>a, Z'>a) = L(a, a; -\rho)
    \end{equation}
    This is the cumulative distribution function (CDF) of the bivariate normal distribution.

    \item \textbf{Term $I_1$ (First-order moment)}:
    \begin{equation}
        I_1 = (1+\rho)\phi(a)\Phi\left(a\sqrt{\frac{1-\rho}{1+\rho}}\right)
    \end{equation}
    where $\phi$ and $\Phi$ are the PDF and CDF of the standard normal distribution.

    \item \textbf{Term $I_2$ (Cross second-order moment)}:
    \begin{equation}
        I_2 = \rho L(a, a; -\rho) + (1-\rho^2)\phi(a,a;\rho)
    \end{equation}
     where $\phi(a,a;\rho) = \frac{1}{2\pi\sqrt{1-\rho^2}}\exp\left(-\frac{a^2}{1+\rho}\right)$.
\end{itemize}

\section{Multi-component and Multi-layer Decomposition}
The effectiveness of MMNNs stems from their ability to decompose a complex function into a combination of smoother, more manageable components. This "divide and conquer" strategy is formalized below, first in one dimension and then generalized to higher dimensions.

\subsection{One-Dimensional Decomposition}
Given a function $f$ on an interval $[x_0, x_n]$ partitioned by points $x_0 < x_1 < \cdots < x_n$, we can decompose $f$ using piecewise linear transformations. Let $\calL_i:[a_i,b_i]\to [x_{i-1},x_i]$ be a linear map. We define a smoother version of $f$ on each subinterval as $f_i = f \circ \calL_i$.

The transformation $\psi_i(x)$ maps the subinterval $[x_{i-1}, x_i]$ to $[a_i, b_i]$ and is constant outside. It can be constructed with ReLU functions:
\begin{equation}
    \psi_i(x)=s_i\cdot \text{ReLU}\left( {x-x_{i-1}} \right)
    -s_i\cdot \text{ReLU}\left( {x-x_{i}} \right)+a_i.
\end{equation}
where $s_i = (b_i-a_i)/(x_i-x_{i-1})$.

\begin{theorem}[1D Decomposition]
The target function $f:[x_0,x_n]\to\mathbb{R}$ can be decomposed as:
\begin{equation}
    f(x) = \sum_{i=1}^n f_i\circ \psi_i(x) - \sum_{i=1}^{n-1}f(x_i)
\end{equation}
for any $x \in [x_0, x_n]$.
\end{theorem}
This theorem shows that a complex function can be represented by a sum of compositions of smoother functions ($f_i$) and simple piecewise linear maps ($\psi_i$), which are easily representable by a shallow network.

\subsection{N-Dimensional Decomposition}
The decomposition strategy extends to multiple dimensions by applying the 1D decomposition iteratively for each dimension. For a 2D function $f(x,y)$ on a domain $[x_0,x_n]\times [y_0,y_m]$, the decomposition formula is derived from the principle of inclusion-exclusion (also known as Poincaré's formula), which provides a way to express the function in terms of its values on sub-domains and at boundary points. This leads to the following theorem.

\begin{theorem}[2D Decomposition]
A function $f:[x_0,x_n]\times [y_0, y_m]\to\mathbb{R}$ can be expressed as:
\begin{equation}
    \begin{split}          
        f( x, y) & = \sum_{i=1}^n \sum_{j=1}^m f_{i,j}\Big( \psi_i(x), \phi_j(y)\Big) - \sum_{i=1}^n\sum_{j=1}^{m-1} f_{i,0}\Big(\psi_i(x), y_{j}\Big) \\
        & \quad - \sum_{i=1}^{n-1} \sum_{j=1}^m f_{0,j}\Big( x_i, \phi_j(y)\Big) + \sum_{i=1}^{n-1}\sum_{j=1}^{m-1} f (x_i, y_{j})
    \end{split}
\end{equation}
where $f_{i,j}$, $f_{i,0}$, $f_{0,j}$ are rescaled versions of $f$, and $\psi_i, \phi_j$ are the 1D piecewise linear transformations for each dimension.
\end{theorem}
This formula, based on the inclusion-exclusion principle, generalizes to any number of dimensions, as detailed in the following theorem. The proof of this general theorem can be established by recursion on the dimension $d$.

\begin{theorem}[N-Dimensional Decomposition]
The decomposition for a function $f: \prod_{k=1}^d [x_{k,0}, x_{k, n_k}] \to \mathbb{R}$ on a hyperrectangle is given by:
\begin{equation}
    f(\bm{x}) = \sum_{S \subseteq \{1, \dots, d\}} (-1)^{d-|S|} \sum_{\bm{j}_{\bar{S}}} \left( \mathcal{D}_S \left[ f(\cdot, \bm{x}_{\bar{S}, \bm{j}_{\bar{S}}}) \right] (\bm{x}_S) \right)
\end{equation}
where:
\begin{itemize}
    \item $\bm{x}_S = (x_k)_{k \in S}$ are the variables for dimensions in the subset $S \subseteq \{1, \dots, d\}$.
    \item $\bm{x}_{\bar{S}, \bm{j}_{\bar{S}}} = (x_{k, j_k})_{k \in \bar{S}}$ are grid points for dimensions in the complement $\bar{S}$.
    \item $\mathcal{D}_S$ is the decomposition operator applied to variables in $S$, which rescales the function on each subdomain of the partition, i.e., $\mathcal{D}_S[g](\bm{y}_S) = \sum_{\bm{i}_S} g_{\bm{i}_S}((\psi_{k, i_k}(y_k))_{k \in S})$.
\end{itemize}
This provides a mathematical foundation for the ability of MMNNs to handle high-dimensional, complex functions by breaking them into simpler parts.
\end{theorem}

\subsection{Further Considerations on Scope}
While the MMNN architecture can be extended to include residual connections (ResMMNNs), the corresponding NTK calculations become significantly more complex. Therefore, this analysis will focus on the standard MMNN architecture, with ResMMNNs being a subject for future work.

Furthermore, our analysis relies on a specific parameter initialization strategy tailored to the MMNN structure and the infinite-width regime ($n \to \infty$):
\begin{itemize}
    \item The inner weights $\tilde{\bm{W}}$, which define the basis functions, are drawn from a fixed, order-one distribution (i.e., their variance does not scale with $n$). This is a crucial and non-standard assumption. The rationale is that these weights define a partition of the input space. To achieve a progressively finer discretization of the domain as the width $n$ grows, the basis functions must be distributed across the entire domain, which requires their parameters to be of order one.
    \item The outer, trainable weights $\tilde{\bm{A}}$, which determine the coefficients of the basis functions, are scaled by $1/\sqrt{n}$ (i.e., initialized with a variance of $1/n$). This normalization is essential for the NTK theory: it ensures that the network output, being a sum over $n$ components, converges to a well-behaved Gaussian process in the infinite-width limit. This allows for a stable approximation of the function's values on the discretized domain created by $\tilde{\bm{W}}$.
\end{itemize}
To begin the analysis, we will focus on the Rectified Linear Unit (ReLU) as the activation function. Its property of being positive-homogeneous (i.e., $\text{ReLU}(\alpha x) = \alpha \text{ReLU}(x)$ for $\alpha > 0$) significantly simplifies the kernel computations, making it an ideal choice for the initial derivation.

\section{NTK Computation for a Single-Layer MMFN}
We now derive the Neural Tangent Kernel for a single-layer MMFN. For simplicity, we consider a single scalar output, such that the network is defined as:
\begin{equation}
    h(\bm{x}) = \frac{1}{\sqrt{n}} \sum_{j=1}^n A_j \sigma(\bm{w}_j^\top \bm{x} + b_j) + c
\end{equation}
The trainable parameters are $\theta = \{A_j\}_{j=1}^n \cup \{c\}$. The parameters $\{\bm{w}_j, b_j\}$ are fixed and randomly drawn, with $\bm{w}_j \sim \mathcal{N}(0, \Sigma_w)$ and $b_j \sim \mathcal{N}(\mu_b, \sigma_b^2)$. For now, we will assume $\Sigma_w = I_d$ and $\mu_b = 0$, $\sigma_b^2 = 1$ for simplicity, which we will generalize shortly.

\subsection{From Gradient Products to an Expectation}
The NTK is defined as the dot product of the gradients of the network output with respect to the trainable parameters. The gradient with respect to a single coefficient $A_j$ and the bias $c$ are:
\begin{align}
    \nabla_{A_j} h(\bm{x}) &= \frac{1}{\sqrt{n}} \sigma(\bm{w}_j^\top \bm{x} + b_j) \\
    \nabla_{c} h(\bm{x}) &= 1
\end{align}
The kernel is the sum of the products of these gradients over all trainable parameters:
\begin{equation}
    K(\bm{x}, \bm{x}') = \sum_{j=1}^n \left( \nabla_{A_j} h(\bm{x}) \nabla_{A_j} h(\bm{x}') \right) + \nabla_{c} h(\bm{x}) \nabla_{c} h(\bm{x}') = \frac{1}{n} \sum_{j=1}^n \sigma(\bm{w}_j^\top \bm{x} + b_j)\sigma(\bm{w}_j^\top \bm{x}' + b_j) + 1
\end{equation}
In the infinite-width limit ($n \to \infty$), by the Law of Large Numbers, this sum converges to an expectation over the distribution of the random parameters $(\bm{w}, b)$:
\begin{equation}
    K_{NTK}(\bm{x}, \bm{x}') = \mathbb{E}_{\bm{w},b} \left[ \sigma(\bm{w}^\top \bm{x} + b)\sigma(\bm{w}^\top \bm{x}' + b) \right] + 1
\end{equation}

\subsection{Bivariate Gaussian Framework and the Non-Centered Case}
The pre-activations $g_1 = \bm{w}^\top \bm{x} + b$ and $g_2 = \bm{w}^\top \bm{x}' + b$ are linear transformations of Gaussian random variables, and are thus themselves jointly Gaussian. In the general case, where the bias $b$ may have a non-zero mean, $(g_1, g_2)$ follow a non-centered bivariate Gaussian distribution. Let's denote their parameters as:
\begin{itemize}
    \item Means: $\mu_1 = \mathbb{E}[g_1]$ and $\mu_2 = \mathbb{E}[g_2]$.
    \item Variances: $\sigma_1^2 = \text{Var}(g_1)$ and $\sigma_2^2 = \text{Var}(g_2)$.
    \item Correlation: $\rho = \frac{\text{Cov}(g_1, g_2)}{\sigma_1 \sigma_2}$.
\end{itemize}
The problem is reduced to computing $\mathbb{E}[\sigma(g_1)\sigma(g_2)]$.

A direct calculation of this integral is remarkably complex. The ReLU activations confine the integration domain to a conical region in the high-dimensional space of the parameters $(\bm{w}, b)$. Solving this integral directly would require substantial linear algebra, equivalent to integrating over a dual cone in the space of dual variables. Fortunately, by standardizing the variables and exploiting the properties of the bivariate Gaussian distribution, we arrive at an analytical solution.

\subsection{Solving the Non-Centered Expectation}
Let $(g_1, g_2)$ be non-centered bivariate Gaussians as described above. We standardize them as $g_1 = \sigma_1 Z_1 + \mu_1$ and $g_2 = \sigma_2 Z_2 + \mu_2$, where $(Z_1, Z_2)$ is a standard bivariate normal with correlation $\rho$. The expectation $\mathbb{E}[\text{ReLU}(g_1)\text{ReLU}(g_2)]$ decomposes into three terms:
\begin{equation}
    \mathbb{E}[\sigma(g_1)\sigma(g_2)] = \sigma_1\sigma_2 \cdot I_2 + (\mu_1\sigma_2 + \mu_2\sigma_1) \cdot I_1 + \mu_1\mu_2 \cdot I_0
\end{equation}
where we have simplified the symmetric case $\mathbb{E}[Z_1|\dots]=\mathbb{E}[Z_2|\dots]$ for clarity. The terms $I_0, I_1, I_2$ are defined for standardized variables $Z_1, Z_2$ with thresholds $a_1 = -\mu_1/\sigma_1$ and $a_2 = -\mu_2/\sigma_2$. For simplicity, we present the result for $a_1=a_2=a$:
\begin{itemize}
    \item \textbf{Term $I_0$ (Probability Term):} This is the probability that both variables are active.
    \begin{equation}
        I_0 = P(Z_1 > a, Z_2 > a) = L(-a, -a; \rho)
    \end{equation}
    where $L$ is the standard bivariate normal CDF.
    \item \textbf{Term $I_1$ (First-Order Moment):}
    \begin{equation}
        I_1 = \mathbb{E}[Z_1 \cdot \mathbf{1}_{Z_1>a, Z_2>a}] = (1+\rho)\phi(a)\Phi\left(a\sqrt{\frac{1-\rho}{1+\rho}}\right)
    \end{equation}
    where $\phi$ and $\Phi$ are the standard normal PDF and CDF, respectively.
    \item \textbf{Term $I_2$ (Second-Order Cross-Moment):}
    \begin{equation}
        I_2 = \mathbb{E}[Z_1 Z_2 \cdot \mathbf{1}_{Z_1>a, Z_2>a}] = \rho \cdot L(-a, -a; -\rho) + \phi(a, a; \rho)
    \end{equation}
    where $\phi(a,a;\rho)$ is the bivariate PDF: $\phi(a,a;\rho) = \frac{1}{2\pi\sqrt{1-\rho^2}} \exp\left(-\frac{a^2}{1+\rho}\right)$.
\end{itemize}
These components provide the complete solution for the expectation, which in turn defines the NTK for the single-layer MMFN.

\section{NTK for a Multi-Layer Network}
\textit{This section will be developed later.}

\section{Additional Conceptual Notes}
The documents contain several conceptual notes on strategy and interpretation.
\begin{itemize}
    \item \textbf{Decomposition:} An idea of decomposition via the inclusion-exclusion principle is mentioned.
    \item \textbf{Scale Interaction:} The handwritten notes mention a "Particle-large-scale-interaction summation".
    \item \textbf{Data Signature:} A note suggests a link between the fineness of information in the data, NTK non-triviality, and the fine "signature" of the result. ("Close A information of the dataset $\rightarrow$ NTK not trivial $\rightarrow$ dx on a zone fine $\rightarrow$ Fine Signature").
    \item \textbf{Homogeneity and \textit{ILLEGIBLE}}: The notion of activation function homogeneity is mentioned in connection with kernel calculation.
\end{itemize}
The final note "Finished $\rightarrow$ NTK Theorem for deep... $\rightarrow$ next paper" suggests that these calculations are a preparatory step for a broader analysis of deep networks.

\end{document}